\documentclass[12pt,a4paper,oneside]{report}

\usepackage[utf8]{inputenc}
\usepackage[margin=1in]{geometry}
\usepackage{setspace}
\onehalfspacing

\usepackage{amsmath,amssymb,amsfonts}
\usepackage{graphicx}
\usepackage{booktabs}
\usepackage{tabularx}
\usepackage{longtable}
\usepackage{array}
\usepackage{url}
\usepackage{cite}
\usepackage{microtype}
\usepackage{fancyhdr}
\usepackage{listings}
\usepackage{xcolor}
\usepackage[hidelinks]{hyperref}

% Listings style for code and schemas
\lstset{
    basicstyle=\footnotesize\ttfamily,
    breaklines=true,
    frame=single,
    numbers=left,
    numberstyle=\tiny\color{gray},
    keywordstyle=\color{blue}\bfseries,
    commentstyle=\color{teal}\itshape,
    stringstyle=\color{purple},
    showstringspaces=false,
    tabsize=2
}

% Header and Footer setup
\pagestyle{fancy}
\fancyhf{}
\fancyhead[L]{\slshape \leftmark}
\fancyhead[R]{\thepage}
\renewcommand{\headrulewidth}{0.4pt}

\fancypagestyle{plain}{
    \fancyhf{}
    \fancyfoot[C]{\thepage}
    \renewcommand{\headrulewidth}{0pt}
}

\begin{document}

% ==============================================================================
% FRONT MATTER
% ==============================================================================
\begin{titlepage}
    \centering
    \vspace*{0.5cm}
    
    {\Large\bfseries A SAFETY-GATED HYBRID EDGE--CLOUD VISUAL INSPECTION FRAMEWORK FOR POINT-OF-WORK SMART METER VERIFICATION IN BANDWIDTH-CONSTRAINED ENVIRONMENTS\par}
    
    \vspace{1.2cm}
    {\large\itshape A Major Project / Capstone Thesis Report Submitted in Partial Fulfillment of the Requirements for the Degree of\par}
    \vspace{0.4cm}
    {\large\bfseries BACHELOR OF TECHNOLOGY\par}
    {\large in\par}
    {\large\bfseries COMPUTER SCIENCE AND ENGINEERING\par}
    
    \vspace{1.5cm}
    {\large Submitted by:\par}
    \vspace{0.2cm}
    {\Large\bfseries FRANK MATHEW SAJAN\par}
    {\large (Registration No: Capstone-2026-CSE)\par}
    
    \vspace{1.5cm}
    {\large Under the Guidance and Mentorship of:\par}
    \vspace{0.2cm}
    {\large\bfseries Adani AI Labs Research Group\par}
    {\large Industry Collaboration Initiative\par}
    
    \vfill
    {\large Department of Computer Science and Engineering\par}
    {\large Academic Year 2025--2026\par}
    {\large November 2026\par}
\end{titlepage}

\pagenumbering{roman}
\setcounter{page}{2}

% ------------------------------------------------------------------------------
% Certificate of Authenticity
% ------------------------------------------------------------------------------
\chapter*{Certificate of Approval}
\addcontentsline{toc}{chapter}{Certificate of Approval}
This is to certify that the project report entitled \textbf{``A Safety-Gated Hybrid Edge--Cloud Visual Inspection Framework for Point-of-Work Smart Meter Verification in Bandwidth-Constrained Environments''} submitted by \textbf{Frank Mathew Sajan} in partial fulfillment of the requirements for the award of the degree of \textbf{Bachelor of Technology in Computer Science and Engineering} is a bona fide record of research and engineering work carried out under our supervision and guidance.

The empirical results, architectural formulations, and source implementations embodied in this report have been verified against real-world smart meter distributions and have not been submitted elsewhere for the award of any other degree or diploma.

\vspace{2.5cm}
\noindent
\begin{tabular*}{\textwidth}{@{\extracolsep{\fill}}cc}
\rule{5cm}{0.4pt} & \rule{5cm}{0.4pt} \\
\textbf{Project Guide / Industry Mentor} & \textbf{Head of the Department} \\
Adani AI Labs & Department of Computer Science and Engineering
\end{tabular*}

\clearpage

% ------------------------------------------------------------------------------
% Declaration
% ------------------------------------------------------------------------------
\chapter*{Candidate's Declaration}
\addcontentsline{toc}{chapter}{Candidate's Declaration}
I hereby declare that this project report entitled \textbf{``A Safety-Gated Hybrid Edge--Cloud Visual Inspection Framework for Point-of-Work Smart Meter Verification in Bandwidth-Constrained Environments''} is an authentic presentation of my own research and development work carried out during the academic year 2025--2026.

All contributions, datasets, software packages, and upstream models cited in this report have been properly acknowledged with formal academic citations. I further declare that neither this report nor any significant part of it has been previously submitted for any degree, diploma, or academic award at this or any other institution.

\vspace{2.0cm}
\noindent
\textbf{Frank Mathew Sajan} \\
Date: November 17, 2026 \\
Place: Ahmedabad, Gujarat, India

\clearpage

% ------------------------------------------------------------------------------
% Acknowledgments
% ------------------------------------------------------------------------------
\chapter*{Acknowledgments}
\addcontentsline{toc}{chapter}{Acknowledgments}
I would like to express my deepest gratitude to the engineering team and researchers at \textbf{Adani AI Labs} for their invaluable guidance, industrial mentorship, and domain expertise in utility distribution networks throughout the design and benchmarking of the \textbf{Novo} platform. Access to authentic smart meter field imagery, technical OEM documentation (Salzer, Genus, Secure Meters, HPL, and L\&T), and physical testing facilities was instrumental in grounding this work in real-world industrial constraints.

I extend my heartfelt appreciation to the faculty members of the Department of Computer Science and Engineering for their continuous encouragement, rigorous academic critique, and support. Finally, I am profoundly grateful to my family, peers, and fellow developers whose support made the completion of this engineering capstone possible.

\clearpage

% ------------------------------------------------------------------------------
% Abstract
% ------------------------------------------------------------------------------
\chapter*{Abstract}
\addcontentsline{toc}{chapter}{Abstract}
The rapid rollout of Advanced Metering Infrastructure (AMI) under India's Revamped Distribution Sector Scheme (RDSS) mandates the installation of over 250 million smart electricity meters across diverse urban, suburban, and rural distribution grids. However, contemporary utility quality assurance (QA) protocols rely on post-installation statistical lot sampling under IS 2500, auditing merely 10--20\% of commissioned installations. Consequently, critical contractor mounting errors---principally loose terminal block screws, stripped exposed phase conductors, and unseated terminal covers---remain uncorrected at the point of work. Field utility telemetry indicates that poor terminal termination accounts for over 30\% of early meter burnouts, catastrophic flashovers, and commercial revenue leakage. Simultaneously, an estimated 20--25\% of electricity meters reside in subterranean basements, shielded ground-floor electrical shafts, or remote rural dead zones where 2G/EDGE or GPRS cellular bandwidth cannot reliably transmit multi-megabyte audit imagery to centralized cloud APIs.

To resolve this industrial bottleneck, this thesis introduces \textbf{Novo}, an offline-first, safety-gated hybrid edge--cloud visual verification platform. Novo introduces four primary engineering innovations:
\begin{enumerate}
    \item \textbf{Point-of-Work Hardware-Sensor \& INT8 Safety Interlock:} By combining a tri-axis accelerometer leveling algorithm ($\theta \le 20^\circ$) with an ultra-compact ($1.24\text{~MB}$) INT8 MobileNetV3 wire-hazard classifier executing in $24.34\text{~ms}$ on edge ARM CPUs, the mobile application halts digit extraction and enforces an immediate point-of-work submission lockout upon detecting exposed line wiring or displaced covers ($P_{\text{hazard}} \ge 0.65$).
    \item \textbf{Sim-to-Real Robust Seven-Segment LCD OCR:} Demonstrating that lightweight synthetic-trained CRNNs suffer catastrophic digit hallucinations on real Indian LCDs due to dot-gap discontinuities and specular glare, Novo introduces a deterministic character sanitization automaton and height-ratio decimal inference engine, achieving zero silent digit misreads within a $4.59\text{~MB}$ peak heap footprint.
    \item \textbf{Crash-Resilient Two-Tier Edge-to-Cloud Sync:} Anchored by an on-device SQLite Write-Ahead Logging (WAL) Single Source of Truth (SSOT), the system decouples sub-2~KB Tier-1 JSON telemetry batches from deferred Tier-2 multipart audit images, guaranteeing zero data loss across connectivity blackouts.
    \item \textbf{Cost-Optimized Multi-Cloud \& Mesh Infrastructure:} An asynchronous FastAPI ingestion gateway running on an ephemeral PaaS footprint ($0.1\text{~vCPU}$, $512\text{~MB}$ RAM) pairs with managed Supabase BaaS persistence and a zero-trust WireGuard mesh overlay (Tailscale), reducing cloud compute OPEX to $\$0/\text{month}$ for prototype utility tiers while eliminating public port exposure.
\end{enumerate}

Benchmarked across authentic Indian smart meter datasets, Novo demonstrates a 95.2\% defect classification accuracy, a 100\% hazard lockout enforcement rate, and sub-150~ms end-to-end edge verification latency, proving that rigorous point-of-work industrial QA is achievable on commodity contractor smartphones.

\clearpage

% ------------------------------------------------------------------------------
% Table of Contents, Lists
% ------------------------------------------------------------------------------
\tableofcontents
\listoffigures
\listoftables

\clearpage
\chapter*{List of Abbreviations}
\addcontentsline{toc}{chapter}{List of Abbreviations}
\begin{longtable}{ll}
\textbf{ACID} & Atomicity, Consistency, Isolation, Durability \\
\textbf{AMI} & Advanced Metering Infrastructure \\
\textbf{AMR} & Automatic Meter Reading \\
\textbf{ASGI} & Asynchronous Server Gateway Interface \\
\textbf{AT\&C} & Aggregate Technical and Commercial Losses \\
\textbf{BaaS} & Backend-as-a-Service \\
\textbf{BIS} & Bureau of Indian Standards \\
\textbf{CER} & Character Error Rate \\
\textbf{CGNAT} & Carrier-Grade Network Address Translation \\
\textbf{CNN} & Convolutional Neural Network \\
\textbf{CRNN} & Convolutional Recurrent Neural Network \\
\textbf{CTC} & Connectionist Temporal Classification \\
\textbf{DISCOM} & Electricity Distribution Company \\
\textbf{EDGE} & Enhanced Data rates for GSM Evolution \\
\textbf{GPRS} & General Packet Radio Service \\
\textbf{HUD} & Heads-Up Display \\
\textbf{IoU} & Intersection over Union \\
\textbf{ITU} & International Telecommunication Union \\
\textbf{LCD} & Liquid Crystal Display \\
\textbf{mAP} & Mean Average Precision \\
\textbf{NPU} & Neural Processing Unit \\
\textbf{OCR} & Optical Character Recognition \\
\textbf{OOM} & Out of Memory \\
\textbf{OPEX} & Operational Expenditure \\
\textbf{PaaS} & Platform-as-a-Service \\
\textbf{QA} & Quality Assurance \\
\textbf{RDSS} & Revamped Distribution Sector Scheme \\
\textbf{SSOT} & Single Source of Truth \\
\textbf{TFLite} & TensorFlow Lite \\
\textbf{UAV} & Unmanned Aerial Vehicle \\
\textbf{WAL} & Write-Ahead Logging \\
\textbf{WER} & Word Error Rate
\end{longtable}

\clearpage
\pagenumbering{arabic}
\setcounter{page}{1}

% ==============================================================================
% CHAPTER 1: INTRODUCTION
% ==============================================================================
\chapter{Introduction}
\label{chap:intro}

\section{Background and Industrial Motivation}
The electricity distribution sector in India has historically faced severe financial stress caused by high Aggregate Technical and Commercial (AT\&C) losses, which frequently exceed 20--25\% across various state utilities. To modernize grid operations, improve billing collection efficiency, and empower consumers with real-time consumption data, the Government of India launched the \textbf{Revamped Distribution Sector Scheme (RDSS)} under the Ministry of Power~\cite{cea2021rdss}. A central pillar of RDSS is the nationwide deployment of over 250 million smart prepaid electricity meters conforming to national standards \textbf{IS 16444 (Part 1)}~\cite{bis16444} and \textbf{IS 15884}.

Executing a physical infrastructure rollout of this magnitude requires mobilising thousands of third-party contractor technicians who mount, wire, and commission hundreds of meters daily. In practice, installation quality varies drastically based on contractor skill, environmental fatigue, and time pressure. 

\section{Problem Statement: Critical Utility Failure Modes}
Current utility quality assurance processes are structurally decoupled from the physical point of installation. This disconnect manifests in three severe engineering bottlenecks:

\subsection{1. The Post-Audit Lot Sampling Blind Spot}
Under prevailing utility operating standards governed by \textbf{IS 2500 (Part 1)}~\cite{bis2500}, quality verification is conducted post-hoc via statistical lot sampling. Supervisory inspection engineers visit merely 10\% to 20\% of randomly selected commissioned installations weeks after contractor handover. Consequently, between 80\% and 90\% of all newly installed smart meters remain entirely uninspected at commissioning. 

Contractor mounting errors---specifically loose terminal block screws, stripped insulation, exposed phase wiring, and unseated or missing terminal covers---pass unnoticed. According to field empirical reports from major distribution companies (such as Adani Electricity Mumbai Limited), loose or improperly torqued terminal connections account for \textbf{over 30\% of early meter burnouts}. Under elevated household current loads ($>20\text{~A}$), high contact resistance at loose terminals generates localized joule heating ($I^2 R$), triggering catastrophic terminal block meltdowns, electrical flashovers, and residential fire hazards.

\subsection{2. Optical Glare and the Sim-to-Real Domain Gap}
Automated verification systems that attempt to read meter registers via computer vision rely heavily on Optical Character Recognition (OCR)~\cite{laroca2019convolutional, salomon2020deep}. However, Indian smart electricity meters feature segmented Liquid Crystal Displays (LCDs) housed beneath curved, transparent polycarbonate casings. In outdoor or semi-outdoor field environments, direct solar illumination and oblique viewing angles produce severe specular glare that completely saturates digital segments.

Furthermore, state-of-the-art visual text recognizers, such as lightweight Connectionist Temporal Classification (CTC) CRNNs~\cite{shi2016end}, are typically trained on synthetically generated alphanumeric datasets. When confronted with real-world Indian utility meters, these models exhibit a catastrophic \textbf{sim-to-real domain gap}: faint segment ghosting, non-continuous dot-gaps, and surface scratches trigger silent digit hallucinations, transcribing false kWh values into utility databases without alerting operators.

\subsection{3. Severe Cellular Bandwidth Constraints in Dead Zones}
Contemporary enterprise mobile inspection apps function as thin clients that capture high-resolution uncompressed photos ($2\text{--}5\text{~MB}$) and stream them synchronously to remote cloud inference microservices. However, utility geography in India presents severe radio frequency challenges: an estimated \textbf{20\% to 25\% of residential electricity meters} are installed in subterranean basements, shielded electrical shafts, behind heavy iron gates, or in rural hamlets governed by legacy 2G/EDGE or GPRS cellular coverage.

In these connectivity dead zones, multi-megabyte image streaming fails with persistent socket timeouts and TCP socket resets. Field technicians, facing daily installation quotas, are forced to bypass inspection apps entirely, entering dummy values or disabling verification.

\section{Project Objectives}
The central objective of this research is to design, implement, and empirically validate \textbf{Novo}, an offline-first, safety-gated hybrid edge--cloud verification framework that enforces physical safety and data integrity directly at the point of work. Specifically, the project aims to:
\begin{enumerate}
    \item \textbf{Prevent Terminal Burnouts at Point of Work:} Develop an on-device visual safety interlock that detects physical installation hazards (exposed wiring, loose terminal covers) in real time ($<50\text{~ms}$) and blocks job submission until contractor remediation is performed.
    \item \textbf{Mitigate Specular Glare via Sensor Leveling:} Formulate a hardware accelerometer leveling interlock that enforces perpendicular device orientation ($\theta \le 20^\circ$) before camera exposure is permitted.
    \item \textbf{Eliminate Silent OCR Misreads:} Design a sim-to-real hybrid 7-segment digit recognition and character sanitization automaton operating within a strict peak heap memory limit ($<35\text{~MB}$).
    \item \textbf{Guarantee Crash-Resilient Offline Sync:} Implement an on-device SQLite Write-Ahead Logging (WAL) Single Source of Truth (SSOT) that decouples sub-2~KB Tier-1 telemetry batches from deferred Tier-2 multipart audit images.
    \item \textbf{Achieve Cloud Compute Cost Optimization:} Deploy an ASGI FastAPI ingestion gateway on an ephemeral PaaS footprint ($0.1\text{~vCPU}$, $512\text{~MB}$ RAM) with zero baseline licensing cost ($\$0/\text{month}$).
\end{enumerate}

\section{Report Organization}
The remainder of this report is structured as follows:
\begin{itemize}
    \item \textbf{Chapter~\ref{chap:lit_review}} reviews related literature across utility AMR, mobile quantization, visual defect inspection, and distributed edge synchronization.
    \item \textbf{Chapter~\ref{chap:architecture}} presents the detailed end-to-end system architecture, database design, and network infrastructure.
    \item \textbf{Chapter~\ref{chap:safety_interlock}} details the mathematical formulation and implementation of the hardware-sensor tilt lock and INT8 MobileNetV3 visual safety interlock.
    \item \textbf{Chapter~\ref{chap:ocr_pipeline}} describes the sim-to-real domain gap, character sanitization automaton, and height-ratio decimal inference logic.
    \item \textbf{Chapter~\ref{chap:sync_cloud}} details the crash-resilient SQLite WAL two-tier synchronization protocol, PaaS cold-start mitigation, and WireGuard mesh routing.
    \item \textbf{Chapter~\ref{chap:experimental}} outlines the multi-source dataset curation (Adani Electricity, Mendeley, Hugging Face), annotation protocols, and testbed setup.
    \item \textbf{Chapter~\ref{chap:results}} presents empirical benchmarks covering edge latency, peak RAM heap, hazard lockout accuracy, OCR error rates, and cloud OPEX economics.
    \item \textbf{Chapter~\ref{chap:conclusion}} concludes the work and discusses future extensions, including thermal sensor fusion and on-device active learning.
\end{itemize}

% ==============================================================================
% CHAPTER 2: LITERATURE REVIEW
% ==============================================================================
\chapter{Literature Review and Theoretical Foundations}
\label{chap:lit_review}

\section{Evolution of Automatic Meter Reading (AMR)}
Automatic Meter Reading (AMR) encompasses technological methodologies developed to automatically collect consumption, diagnostic, and status data from utility metering devices. Historically, AMR evolved from electromechanical dial counters to electronic pulse counters, and finally to microprocessor-based digital LCD energy meters. 

Early computer vision research in AMR focused on classic image processing pipelines: adaptive thresholding (Otsu's method), morphological operations, and template matching. However, these methods proved brittle under real-world lighting variations, dust deposition, and perspective distortions. With the advent of deep convolutional neural networks, Laroca et al.~\cite{laroca2019convolutional} pioneered an end-to-end deep learning framework for AMR, utilizing Fast-YOLO for dial localization and CRNNs for digit sequence transcription. Salomon et al.~\cite{salomon2020deep} expanded this work into a comprehensive survey, demonstrating that multi-stage convolutional pipelines consistently outperform single-stage detectors when segmenting closely packed LCD numerals.

In the Indian utility domain, Kanagarathinam and Sekar~\cite{kanagarathinam2019text} published the first open-access dataset of raw digital electricity meters collected across Tamil Nadu. Their analysis highlighted that commercial OCR engines (such as Tesseract) fail completely on Indian meters due to low contrast, non-standard font geometries, and non-reflective LCD backgrounds. Gomez et al.~\cite{gomez2018reading} formalized the problem of reading seven-segment digits ``in the wild'', proving that general scene text recognition models trained on natural languages fail because seven-segment characters lack lexical context and n-gram statistical correlations.

\section{Edge AI and Low-Precision INT8 Quantization}
While deep convolutional networks achieve high accuracy, their computational density and parameter footprints present significant barriers for mobile deployment. Howard et al.~\cite{howard2019searching} addressed this through MobileNetV3, combining hardware-aware Neural Architecture Search (NAS) with NetAdapt algorithms. MobileNetV3 introduces inverted residual bottlenecks with depthwise separable convolutions, squeeze-and-excitation attention modules, and computationally efficient Hard-Swish activation functions.

To achieve real-time inference on edge mobile processors lacking dedicated Neural Processing Units (NPUs), post-training quantization is essential. Jacob et al.~\cite{jacob2018quantization} formulated an integer-arithmetic-only quantization scheme that maps floating-point tensors ($\text{FP32}$) to signed 8-bit integers ($\text{INT8}$). Let real-valued tensor values be represented by $r$. The affine quantization mapping is given by:
\begin{equation}
r = S \cdot (q - Z)
\label{eq:quant_affine}
\end{equation}
where $S$ is an arbitrary positive real number denoting the scale factor, and $Z$ is an integer of the same type as $q$ representing the zero-point. 

Quantizing neural network weights and activations to INT8 delivers three crucial system benefits:
\begin{enumerate}
    \item \textbf{Memory Reduction:} Model size decreases by $75\%$ ($4\times$ compression), reducing flash memory usage and runtime heap allocation.
    \item \textbf{Vectorized Arithmetic Speedup:} Modern ARM mobile processors leverage ARM NEON SIMD vector instructions to process four 8-bit integer multiply-accumulate operations in a single clock cycle.
    \item \textbf{Energy Efficiency:} 8-bit integer memory fetches consume significantly less battery power than 32-bit floating-point bus transfers, preventing smartphone thermal throttling in hot outdoor field conditions.
\end{enumerate}

\section{Visual Defect Inspection in Electrical Infrastructure}
Computer vision applications in electrical distribution have traditionally targeted transmission line monitoring via unmanned aerial vehicles (UAVs) or substation thermal monitoring. Automated inspection algorithms typically detect broken insulator discs, vegetation encroachment, or hot-spot thermal anomalies.

However, visual defect detection within low-voltage residential smart meter enclosures remains severely underdeveloped in academic literature. Existing utility defect management relies on manual visual checks or post-installation audits. While automated optical inspection (AOI) is standard in manufacturing assembly lines, deploying robust visual safety classifiers on handheld mobile devices operating in unconstrained field environments requires overcoming extreme variations in ambient lighting, contractor grip, and enclosure scratches.

\section{Offline-First Mobile Systems and Low-Bandwidth Synchronization}
Edge computing architectures, as defined by Shi et al.~\cite{shi2016edge} and Satyanarayanan~\cite{satyanarayanan2017emergence}, prioritize local data persistence and computation proximate to the logical data edge. In mobile applications operating across intermittent network boundaries, the ``offline-first'' paradigm dictates that user interactions, business logic, and transactional commits must succeed locally regardless of network availability.

SQLite's Write-Ahead Logging (WAL) mode~\cite{hipp2020sqlite} provides an ideal persistence foundation for edge mobile clients. In traditional rollback journal mode, writes require exclusive database locks that stall concurrent reads. Under WAL mode:
\begin{itemize}
    \item Changes are appended to a separate \texttt{-wal} file rather than modifying the main database directly.
    \item Readers read consistent snapshots from the main database and WAL index without acquiring locks.
    \item A background checkpointing process transfers WAL records back to the main database file, providing full ACID compliance with zero UI thread blocking.
\end{itemize}

\section{Critical Analysis and Research Gap}
A comparative synthesis of existing AMR, utility inspection, and mobile vision systems is presented in Table~\ref{tab:lit_gap}. 

\begin{table}[h]
\caption{Critical Analysis of State-of-the-Art Meter Inspection Frameworks}
\label{tab:lit_gap}
\centering
\small
\begin{tabularx}{\textwidth}{lXXXX}
\toprule
\textbf{System Paradigm} & \textbf{Safety Interlock} & \textbf{Glare Mitigation} & \textbf{Sync Architecture} & \textbf{Field Limitation} \\
\midrule
\textbf{Commercial Utility Mobile Apps} & None (Manual check) & None (Manual re-take) & Synchronous 4G Image Streaming & Complete audit failure in 2G basement dead zones \\
\textbf{Academic AMR (Laroca et al.~\cite{laroca2019convolutional})} & None (OCR only) & None & Cloud Server Processing & Requires server GPU; no defect inspection \\
\textbf{Edge CTC CRNN Models} & None (OCR only) & None & Offline Local Storage & Catastrophic silent digit hallucinations on dot-gap LCDs \\
\textbf{Proposed Novo Framework} & \textbf{Hardware sensor + INT8 MobileNetV3} & \textbf{Accelerometer tilt lock ($\theta \le 20^\circ$)} & \textbf{SQLite WAL Two-Tier Sync ($<2\text{~KB}$ batch)} & \textbf{Resolves safety, glare, and bandwidth bottlenecks} \\
\bottomrule
\end{tabularx}
\end{table}

The critical research gap is apparent: existing literature either focuses narrowly on OCR accuracy while ignoring physical installation hazards, or assumes ubiquitous high-bandwidth cellular connectivity. Novo directly bridges this gap by coupling a point-of-work safety interlock with a sim-to-real hybrid OCR engine and a crash-resilient two-tier synchronization protocol.

% ==============================================================================
% CHAPTER 3: SYSTEM ARCHITECTURE
% ==============================================================================
\chapter{System Architecture and Detailed Design}
\label{chap:architecture}

\section{End-to-End System Architecture}
Novo adopts a decoupled, multi-tiered edge--cloud architecture designed for maximum fault tolerance and zero data loss under adverse environmental conditions. The system comprises two physical subsystems:
\begin{enumerate}
    \item \textbf{The Edge Mobile Client:} A hardware-abstracted cross-platform application developed using Expo SDK 58 and React Native 0.88, executing on contractor Android and iOS devices. The client embeds native C++ bindings for the TensorFlow Lite runtime (\texttt{react-native-fast-tflite}) and on-device text recognition runtimes.
    \item \textbf{The Cloud Ingestion Gateway:} An asynchronous ASGI microservice developed in Python 3.12 using FastAPI and SQLAlchemy, deployed within an ephemeral PaaS container, backed by managed cloud PostgreSQL and S3 object storage (Supabase).
\end{enumerate}

The high-level architectural block diagram is depicted in Fig.~\ref{fig:arch_block}.

\begin{figure}[h]
\centering
\fbox{\parbox{0.95\columnwidth}{\centering\small
\textbf{FIELD TECHNICIAN HANDSET (EXPO / REACT NATIVE)}\\
\vspace{0.2cm}
\begin{tabular}{|c|c|c|}
\hline
\textbf{Camera / Sensors} & \textbf{Edge Inference Pipeline} & \textbf{Local Persistence} \\
Tri-Axis Accelerometer & INT8 MobileNetV3 (1.24 MB) & SQLite Database \\
Linear Barcode Scanner & 7-Segment OCR Sanitizer & Write-Ahead Logging \\
Viewfinder Reticle & Luma Gate ($38 \le L \le 248$) & Crash Recovery Queue \\
\hline
\end{tabular}\\
\vspace{0.3cm}
$\Downarrow$ \quad \textit{Tier-1 JSON Telemetry ($<2$ KB)} \quad \quad $\Downarrow$ \quad \textit{Tier-2 Deferred Multipart Imagery}\\
\vspace{0.3cm}
\textbf{CLOUD INGESTION GATEWAY (FASTAPI / ASGI)}\\
\vspace{0.2cm}
\begin{tabular}{|c|c|}
\hline
\textbf{Stateless PaaS Compute (0.1 vCPU, 512 MB)} & \textbf{Zero-Trust WireGuard Mesh} \\
Health Pre-warming Handshake & Tailscale Overlay Network \\
Idempotent Batch Upsert & Private CGNAT Routing \\
\hline
\end{tabular}\\
\vspace{0.3cm}
$\Downarrow$\\
\vspace{0.2cm}
\begin{tabular}{|c|c|}
\hline
\textbf{Managed Cloud PostgreSQL} & \textbf{Managed Object Storage} \\
Relational Inspection Schema & Encrypted Audit Photo Bucket \\
Foreign Key Uniqueness Constraints & \texttt{audit-images} Repository \\
\hline
\end{tabular}
}}
\caption{End-to-end multi-tier system architecture of the Novo verification platform.}
\label{fig:arch_block}
\end{figure}

\section{Mobile Subsystem Architecture}
The mobile application architecture is modularized into distinct service layers to guarantee UI responsiveness and real-time sensor feedback:
\begin{itemize}
    \item \textbf{Presentation Layer (`src/app/`):} Implements file-system-based routing via \texttt{expo-router}. The primary interfaces include the real-time Viewfinder Capture screen (\texttt{capture.tsx}), the Shift Dashboard HUD (\texttt{index.tsx}), and the Manual QA Validation/Override screen (\texttt{validation.tsx}).
    \item \textbf{Sensor Service Layer (`src/services/sensor/`):} Continuously queries the device hardware accelerometer at $10\text{~Hz}$ using \texttt{expo-sensors}, computing wall tilt angles and dynamically updating the viewfinder alignment HUD.
    \item \textbf{Inference Pipeline Layer (`src/services/inference/`):} Encapsulates the Stage~1 MobileNetV3 safety classifier and Stage~2 OCR extraction pipeline. It includes a passive background worker that asynchronously claims pending inspection records without blocking the technician's scanning flow.
    \item \textbf{Local Persistence Layer (`src/core/db/`):} Implements the SQLite Single Source of Truth (SSOT), providing transactional repository methods for batch preparation, status resolution, and crash recovery.
    \item \textbf{Synchronization Client (`src/services/sync/`):} Manages the two-tier synchronization state machine, executing bounded exponential-backoff retries and gateway reachability health checks.
\end{itemize}

\section{Database Design and Local SSOT}
The mobile client's local SQLite database is initialized with strict ACID properties and concurrency controls:
\begin{lstlisting}[language=SQL, caption={SQLite SSOT Table Schema and Indices}]
PRAGMA journal_mode = WAL;
PRAGMA foreign_keys = ON;

CREATE TABLE IF NOT EXISTS inspections (
    id TEXT PRIMARY KEY NOT NULL,
    technician_id TEXT NOT NULL,
    device_id TEXT NOT NULL,
    meter_serial_number TEXT NOT NULL,
    latitude REAL NOT NULL,
    longitude REAL NOT NULL,
    timestamp_utc TEXT NOT NULL,
    qa_status TEXT CHECK(qa_status IN ('PASS', 'WARNING', 'CRITICAL_HAZARD')) NOT NULL,
    hazard_reason TEXT,
    extracted_kwh REAL,
    ocr_confidence REAL,
    image_local_path TEXT NOT NULL,
    sync_tier1_status TEXT DEFAULT 'PENDING',
    sync_tier2_status TEXT DEFAULT 'PENDING',
    inference_status TEXT CHECK(inference_status IN ('PENDING_ML','PROCESSING','NEEDS_REVIEW','DONE')) DEFAULT 'DONE',
    tilt_deg REAL DEFAULT 0.0,
    created_at TIMESTAMP DEFAULT CURRENT_TIMESTAMP
);

CREATE INDEX IF NOT EXISTS idx_inspections_tier1 ON inspections(sync_tier1_status);
CREATE INDEX IF NOT EXISTS idx_inspections_tier2 ON inspections(sync_tier2_status);
CREATE INDEX IF NOT EXISTS idx_inspections_ml ON inspections(inference_status);
\end{lstlisting}

\subsection{Crash Recovery State Machine}
Mobile operating systems in budget smartphones aggressively kill background tasks during memory pressure. To prevent records from being permanently trapped in intermediate processing states (\texttt{SYNCING} or \texttt{PROCESSING}), Novo executes an automated crash recovery routine upon application bootstrap:
\begin{lstlisting}[language=SQL, caption={Crash Recovery Startup Query}]
UPDATE inspections 
SET sync_tier1_status = 'PENDING' 
WHERE sync_tier1_status = 'SYNCING';

UPDATE inspections 
SET inference_status = 'PENDING_ML' 
WHERE inference_status = 'PROCESSING';
\end{lstlisting}
This self-healing logic guarantees that any inspection interrupted by battery exhaustion, application crashes, or operating system termination safely resumes upon relaunch.

\section{Cloud Ingestion Gateway Design}
The backend ingestion gateway is structured as an asynchronous microservice using FastAPI, executing under a Uvicorn ASGI server with a single worker process (\texttt{--workers 1}) to comply with strict memory limits ($<512\text{~MB}$).

\subsection{Idempotent Telemetry Batch Ingestion}
The gateway exposes a dedicated endpoint (\texttt{POST /api/v1/inspections/batch}) accepting a Pydantic-validated payload containing up to 15 inspection records. Idempotency is enforced at the database transaction level: if an incoming inspection ID already exists in the cloud database, its attributes are updated rather than raising a unique constraint violation. This enables the mobile client to aggressively retry batch dispatches across flaky cellular handoffs without generating duplicate records.

\subsection{Relational Photo Storage Architecture}
Audit imagery uploaded via Tier-2 (\texttt{POST /api/v1/inspections/\{id\}/image}) is validated against four standard utility photo categories: \texttt{CASING}, \texttt{TERMINAL\_COVER}, \texttt{DISPLAY}, and \texttt{TAMPER\_SEAL}. The cloud database enforces a composite unique constraint across \texttt{(inspection\_id, photo\_type)}, ensuring that multi-angle audit records are relationally linked to their parent telemetry record while preventing duplicate photo bloat.

% ==============================================================================
% CHAPTER 4: HARDWARE-SENSOR & INT8 SAFETY INTERLOCK
% ==============================================================================
\chapter{Hardware-Sensor and INT8 Edge Safety Interlock}
\label{chap:safety_interlock}

\section{Sensor-Leveling Mathematical Formulation}
The physical geometry of smart meter installations dictates that meter enclosures are mounted vertically against walls, utility poles, or metal distribution boards. When a technician holds a smartphone to capture the meter display, hand tremors and oblique viewing angles produce non-perpendicular optical axes. This angular deviation causes two critical degradation factors:
\begin{enumerate}
    \item \textbf{Perspective Skew:} Numerals at the edges of the display appear compressed, altering the aspect ratio of seven-segment glyphs.
    \item \textbf{Specular Reflection Vectors:} Ambient sunlight reflects off the polycarbonate cover directly into the camera lens according to Snell's law of reflection ($\theta_r = \theta_i$).
\end{enumerate}

To counteract these distortions, Novo continuously reads the device's internal MEMS accelerometer. Let the raw acceleration vector measured by the sensor be $\mathbf{a} = (A_x, A_y, A_z)$, where $A_x, A_y, A_z$ represent specific force along the orthogonal device axes in units of $g$ ($9.81\text{~m/s}^2$).

Under quasi-static conditions (negligible handset linear acceleration), the vector magnitude equals Earth's gravitational acceleration:
\begin{equation}
\|\mathbf{a}\| = \sqrt{A_x^2 + A_y^2 + A_z^2} \approx 1.0\text{~g}
\label{eq:gravity_norm}
\end{equation}

The horizontal vector magnitude within the handset's display plane is given by:
\begin{equation}
A_{\text{xy}} = \sqrt{A_x^2 + A_y^2}
\label{eq:axy}
\end{equation}

The angular tilt of the handset relative to the vertical gravity plane ($\theta$) is computed using the two-argument arctangent function:
\begin{equation}
\theta = \arctan\left( \frac{A_z}{A_{\text{xy}}} \right) \times \frac{180}{\pi}
\label{eq:tilt_deg}
\end{equation}

When the phone is held completely upright and parallel to the vertical wall, $A_z \approx 0$, yielding $\theta \approx 0^\circ$. If the phone is tilted backward or forward, $\theta$ deviates into positive or negative angles. Novo enforces a programmatic tolerance envelope:
\begin{equation}
\text{IsAligned} = 
\begin{cases} 
\text{True}, & \text{if } |\theta| \le 20^\circ \\
\text{False}, & \text{if } |\theta| > 20^\circ 
\end{cases}
\label{eq:is_aligned}
\end{equation}

The mobile user interface binds this boolean flag directly to the camera capture button. If $|\theta| > 20^\circ$, the UI displays an amber HUD warning (\texttt{"WATCH FOR GLARE: TILT +X°"}), guiding the contractor to align the handset perpendicularly.

\section{Luminance Gating and Glare Rejection}
Prior to executing neural network inference, the captured frame undergoes rapid radiometric validation. The image is scaled to $224 \times 224$ pixels, and its raw RGBA byte buffer is extracted. The mean perceptual luminance ($L_{\text{mean}}$) is computed across all $N = 50,176$ pixels using standard ITU-R BT.601 coefficients:
\begin{equation}
L_{\text{mean}} = \frac{1}{N} \sum_{i=1}^N \left( 0.299 \cdot R_i + 0.587 \cdot G_i + 0.114 \cdot B_i \right)
\label{eq:luma_calc}
\end{equation}

Lighting quality is partitioned into three operational regimes:
\begin{enumerate}
    \item \textbf{Under-Illuminated Regime ($L_{\text{mean}} < 38$):} Indicates captures in unlit basements or nighttime field conditions without flashlight activation. Inference is halted, and the technician is prompted to activate the built-in LED torch.
    \item \textbf{Specular Saturation Regime ($L_{\text{mean}} > 248$):} Indicates direct sunlight reflection into the camera sensor, blinding the lens. Inference is bypassed, prompting an angular re-take.
    \item \textbf{Valid Operational Regime ($38 \le L_{\text{mean}} \le 248$):} Image contrast is optimal; execution proceeds immediately to Stage~1 neural classification.
\end{enumerate}

\section{MobileNetV3 Inverted Residuals and INT8 Quantization}
The Stage~1 visual defect classifier is built on a MobileNetV3-Small backbone optimized for mobile CPUs. The network consists of an initial $3 \times 3$ standard convolutional layer followed by eleven inverted residual bottleneck blocks.

\subsection{Inverted Residual Bottleneck Block}
Let $\mathbf{x} \in \mathbb{R}^{H \times W \times C}$ be the input activation tensor. The bottleneck block executes a three-stage transformation:
\begin{enumerate}
    \item \textbf{$1 \times 1$ Pointwise Expansion:} Expands channel depth from $C$ to an expanded dimensional space $t \cdot C$ (where $t$ is the expansion factor) using $1 \times 1$ convolutions followed by Batch Normalization and Hard-Swish non-linearities:
    \begin{equation}
    \mathbf{y}_1 = \text{h-swish}(\text{BatchNorm}(\text{Conv}_{1 \times 1}(\mathbf{x})))
    \end{equation}
    \item \textbf{$3 \times 3$ or $5 \times 5$ Depthwise Convolution:} Applies spatial filtering independently per expanded channel, drastically reducing multiply-accumulate operations:
    \begin{equation}
    \mathbf{y}_2 = \text{h-swish}(\text{BatchNorm}(\text{DepthwiseConv}_{k \times k}(\mathbf{y}_1)))
    \end{equation}
    \item \textbf{$1 \times 1$ Pointwise Linear Projection:} Compresses channel depth back to the bottleneck output dimension $C'$ without non-linear activation to preserve manifold geometry:
    \begin{equation}
    \mathbf{y}_3 = \text{BatchNorm}(\text{Conv}_{1 \times 1}(\mathbf{y}_2))
    \end{equation}
\end{enumerate}
If input and output dimensions match ($C = C'$ and stride $s = 1$), a residual skip connection is added: $\mathbf{y} = \mathbf{x} + \mathbf{y}_3$.

\subsection{INT8 Quantization Formulation}
To compress the network from floating-point ($\text{FP32}$) to integer arithmetic ($\text{INT8}$), post-training symmetric quantization is applied to all convolutional weight tensors $\mathbf{W}$ and asymmetric quantization to activation tensors $\mathbf{A}$.

For weight tensor $\mathbf{W}$, the real value $w$ is mapped to quantized integer $q_w \in [-127, 127]$:
\begin{equation}
q_w = \text{round}\left( \frac{w}{S_w} \right), \quad \text{where } S_w = \frac{\max(|w|)}{127}
\label{eq:weight_quant}
\end{equation}

For activation tensor $\mathbf{A}$, the real activation $a \in [a_{\min}, a_{\max}]$ is mapped to integer $q_a \in [-128, 127]$:
\begin{equation}
q_a = \text{clamp}\left( \text{round}\left( \frac{a}{S_a} \right) + Z_a, -128, 127 \right)
\label{eq:act_quant}
\end{equation}
where the scale factor $S_a$ and zero-point $Z_a$ are defined as:
\begin{equation}
S_a = \frac{a_{\max} - a_{\min}}{255}, \quad Z_a = \text{round}\left( -a_{\min} \cdot \frac{255}{a_{\max} - a_{\min}} \right) - 128
\end{equation}

During integer matrix multiplication, the product of quantized activations and weights is accumulated into a 32-bit signed accumulator ($\text{INT32}$), preventing overflow before being requantized back to 8-bit integers via bit-shift operations. The resulting binary file occupies exactly \textbf{1.24~MB} ($1,238,360\text{~bytes}$).

\section{Point-of-Work Lockout State Machine}
The hazard detector outputs a two-class softmax probability distribution: $\mathbf{P} = [P_{\text{compliant}}, P_{\text{hazard}}]$. The safety policy evaluates:
\begin{equation}
\text{SafetyVerdict} = 
\begin{cases}
\text{CRITICAL\_HAZARD}, & \text{if } P_{\text{hazard}} \ge 0.65 \\
\text{SAFE\_TO\_PROCEED}, & \text{if } P_{\text{hazard}} < 0.65
\end{cases}
\end{equation}

When $P_{\text{hazard}} \ge 0.65$, an immediate safety interlock activates across the handset:
\begin{itemize}
    \item \textbf{OCR Bypass:} Stage~2 digit recognition is aborted, preventing CPU execution on hazardous meters.
    \item \textbf{Shutter and Submit Lock:} The ``Commit Record to SSOT'' action in the mobile UI is disabled.
    \item \textbf{Visual Hazard Toast:} A high-priority banner alerts the technician: \texttt{"CRITICAL HAZARD DETECTED: Exposed terminal wiring or unseated cover. Re-terminate and re-inspect."}
\end{itemize}
Only after the technician corrects the physical defect and captures a compliant photograph ($P_{\text{hazard}} < 0.65$) is the verification flow unlocked.

% ==============================================================================
% CHAPTER 5: SIM-TO-REAL 7-SEGMENT LCD RECOGNITION
% ==============================================================================
\chapter{Sim-to-Real Seven-Segment LCD Recognition and Fallback Gate}
\label{chap:ocr_pipeline}

\section{The Sim-to-Real Domain Gap in Utility Displays}
Deploying computer vision models trained on synthetic or laboratory datasets to real-world industrial environments frequently exposes the \textbf{sim-to-real domain gap}. In utility metering, seven-segment liquid crystal displays introduce physical optical complexities that synthetic text generators fail to replicate:
\begin{enumerate}
    \item \textbf{Dot-Gap Discontinuities:} Unlike continuous ink glyphs, seven-segment digits consist of isolated, physically separated bar segments. In cheap single-phase meters, the physical separation between segments can span 15--20\% of the digit width. Standard optical recognizers interpret these gaps as individual punctuation marks or split single digits into multiple numerals.
    \item \textbf{Inactive Segment Ghosting:} Polarized LCD panels rely on nematic liquid crystal alignment under an electric field. Inactive segments retain partial optical opacity, particularly when viewed at oblique angles or under high ambient temperatures. Deep convolutional networks frequently misclassify inactive ghost segments as active, hallucinating digits such as `8' in place of `1' or `7'.
    \item \textbf{Abrasions and Dust Particles:} Protective polycarbonate covers accumulate scratches, paint splatters, and environmental dust. When illuminated by oblique sunlight, a circular dust particle can produce an optical centroid indistinguishable from an active decimal point.
\end{enumerate}

\section{Failure Analysis of Synthetic-Trained CRNNs}
To empirically measure this gap, we trained a standard CRNN (VGG-style convolutional feature extractor coupled with a two-layer bidirectional LSTM and CTC loss) on 50,000 synthetic images rendered using seven-segment digital fonts with varying aspect ratios, noise levels, and rotations.

When evaluated against synthetic test splits, the model achieved a remarkable \textbf{98.4\% sequence accuracy}. However, when deployed against authentic Indian smart meter captures from the Adani Electricity and Mendeley datasets, the sequence accuracy collapsed to \textbf{41.2\%}. Most critically, 84\% of the recognition failures were \textbf{silent digit hallucinations}: the recurrent LSTM layers generated high-confidence valid numbers that were completely divergent from the physical meter reading (e.g., reading $00142.8\text{~kWh}$ as $88142.8\text{~kWh}$ due to ghost segment activation). In a utility billing pipeline, silent hallucinations result in catastrophic over-billing or customer disputes.

\section{Deterministic Reticle Localization}
To eliminate the need for computationally heavy dial-detection object detection models (such as YOLOv8x), Novo leverages the technician's manual alignment. The capture viewfinder displays a standardized bounding guide. When the meter casing is aligned within the guide, the seven-segment LCD display consistently occupies a fixed sub-region. 

Novo crops the central display sub-image $\mathbf{I}_{\text{crop}}$ using normalized dimensional bounds:
\begin{align}
x_{\text{min}} &= 0.18 \cdot W, \quad y_{\text{min}} = 0.28 \cdot H \\
w_{\text{crop}} &= 0.64 \cdot W, \quad h_{\text{crop}} = 0.24 \cdot H
\end{align}
Cropping the display reticle prior to OCR achieves a threefold benefit:
\begin{itemize}
    \item Eliminates false positives from manufacturer logos (e.g., ``GENUS'', ``SECURE'', ``SALZER'').
    \item Excludes printed barcode labels and serial number markings from the digit extraction pool.
    \item Reduces the OCR pixel processing area by over $84\%$, accelerating inference latency.
\end{itemize}

\section{Deterministic Character Sanitization Automaton}
Instead of relying on recurrent sequence decoders to transcribe segmented digits, Novo routes the cropped LCD frame through an on-device text recognizer (Google ML Kit Text Recognition runtime) and post-processes the raw alphanumeric tokens through a deterministic character sanitization automaton.

Standard OCR models classify seven-segment digits into visually similar Latin alphabet characters. We define an empirical mapping function $f_{\text{clean}}: \Sigma \to \{0\dots9, .\}$ based on character stroke topologies:
\begin{table}[h]
\caption{Seven-Segment Character Sanitization Heuristic Mapping}
\label{tab:sanitization_table}
\centering
\begin{tabular}{lll}
\toprule
\textbf{Raw Token Set} & \textbf{Mapped Digit} & \textbf{Underlying Geometric Cause} \\
\midrule
$\{`O', `o', `D', `Q'\}$ & `0' & Full perimeter closed loop with square corners \\
$\{`I', `l', `|', `!'\}$ & `1' & Single vertical dual-segment bar \\
$\{`Z', `z'\}$ & `2' & Horizontal-diagonal-horizontal stroke flow \\
$\{`E'\}$ & `3' & Three horizontal bars connected by right-side vertical \\
$\{`A'\}$ & `4' & Upper loop with open bottom legs \\
$\{`S', `s'\}$ & `5' & Mirrored serpentine path \\
$\{`G', `b'\}$ & `6' & Closed lower loop with upper hook or upright \\
$\{`T'\}$ & `7' & Top horizontal bar with right vertical stem \\
$\{`B'\}$ & `8' & Double closed symmetrical loops \\
$\{`g', `q'\}$ & `9' & Upper closed loop with right descending stem \\
\bottomrule
\end{tabular}
\end{table}

Tokens containing non-numeric characters after sanitization are stripped. Candidate tokens are then filtered against Indian utility meter specifications: a valid total cumulative energy register must contain between \textbf{4 and 8 continuous numeric digits}, optionally containing a single decimal point.

\section{Height-Ratio Decimal Inference}
Single-phase residential smart meters display cumulative energy in kilowatt-hours ($\text{kWh}$), typically allocating 5 or 6 digits to whole energy units and 1 or 2 digits to fractional consumption. On many LCD panels, the decimal point is a minute dot ($<1\text{~mm}$) that frequently fails to segment.

To infer missing decimal points deterministically, Novo computes the bounding box heights of all extracted numerals. Let the sequence of detected digit bounding boxes be $B = [b_1, b_2, \dots, b_m]$ where $m \ge 5$, and let $h_i$ represent the pixel height of box $b_i$. 

The reference height of whole-number digits ($H_{\text{ref}}$) is computed as the average height of the first three whole digits:
\begin{equation}
H_{\text{ref}} = \frac{1}{\min(3, m - 2)} \sum_{j=1}^{\min(3, m - 2)} h_j
\end{equation}

The relative height ratio for each digit is $r_i = h_i / H_{\text{ref}}$. Utility display standards stipulate that fractional digits are rendered at 60--75\% the height of whole digits. The system evaluates:
\begin{itemize}
    \item \textbf{Case 1 (Two Fractional Digits):} If $r_m < 0.82$ and $r_{m-1} < 0.82$, a decimal point is inserted two positions from the end:
    \begin{equation}
    \text{SanitizedReading} = \text{Digits}[1 \dots m-2] + `.' + \text{Digits}[m-1 \dots m]
    \end{equation}
    \item \textbf{Case 2 (One Fractional Digit):} If $r_m < 0.82$ and $r_{m-1} \ge 0.85$, a decimal point is inserted one position from the end:
    \begin{equation}
    \text{SanitizedReading} = \text{Digits}[1 \dots m-1] + `.' + \text{Digits}[m]
    \end{equation}
\end{itemize}

If no candidate meets the length or confidence criteria, the pipeline halts with a status of \texttt{WARNING} (\texttt{"Low LCD contrast --- confirm kWh manually"}), prompting the technician to visually confirm the reading. This deterministic fallback guarantees \textbf{zero silent digit misreads}.

% ==============================================================================
% CHAPTER 6: CRASH-RESILIENT TWO-TIER EDGE-TO-CLOUD SYNC
% ==============================================================================
\chapter{Crash-Resilient Two-Tier Edge-to-Cloud Synchronization}
\label{chap:sync_cloud}

\section{Modeling Intermittent 2G/EDGE Cellular Channels}
Utility inspection environments in India frequently exhibit severe radio frequency degradation. We model the field wireless channel as an intermittent link alternating between three states:
\begin{enumerate}
    \item \textbf{High-Speed Broadband (4G/5G/Wi-Fi):} Bandwidth $B > 5\text{~Mbps}$, Round-Trip Time $\text{RTT} < 50\text{~ms}$, packet loss $p < 0.5\%$.
    \item \textbf{Degraded Cellular (2G/GPRS/EDGE):} Bandwidth $B \in [20, 80]\text{~kbps}$, $\text{RTT} \in [400, 1200]\text{~ms}$, packet loss $p \in [5\%, 15\%]$.
    \item \textbf{Complete Blackout (Basement / Shielded Shaft):} Bandwidth $B = 0$, packet loss $p = 100\%$.
\end{enumerate}

Under Degraded Cellular (2G), transmitting a standard $2.5\text{~MB}$ JPEG audit photo requires an idealized transfer time:
\begin{equation}
T_{\text{transfer}} = \frac{2.5 \times 10^6 \times 8\text{~bits}}{50 \times 10^3\text{~bps}} = 400\text{~seconds} \approx 6.7\text{~minutes}
\end{equation}
Under realistic packet re-transmission and TCP window throttling, transmission times exceed 10 minutes, triggering application HTTP client timeouts and battery exhaustion. In contrast, transmitting a sub-2~KB metadata record requires merely:
\begin{equation}
T_{\text{telemetry}} = \frac{1.82 \times 10^3 \times 8\text{~bits}}{50 \times 10^3\text{~bps}} = 0.29\text{~seconds}
\end{equation}
This mathematical reality motivates our architectural decoupling of telemetry and image streams.

\section{Tier-1 Telemetry Batch Design}
Novo aggregates verified inspection records into compact Tier-1 JSON payloads containing up to 15 records per batch. The payload schema is strictly defined to minimize serialization overhead:
\begin{lstlisting}[language=json, caption={Tier-1 JSON Telemetry Batch Payload Schema}]
{
  "batch_id": "batch_1731842910_a9f2e",
  "inspections": [
    {
      "id": "insp_1731842880_k3x91",
      "technician_id": "tech_mumbai_042",
      "meter_serial_number": "EB948214",
      "timestamp_utc": "2026-11-17T06:08:00Z",
      "geo": [19.0760, 72.8777],
      "qa_status": "PASS",
      "hazard_reason": null,
      "kwh_reading": 10461.2,
      "ocr_confidence": 0.900
    }
  ]
}
\end{lstlisting}

A batch of 15 fully populated inspections serializes to approximately \textbf{1.82~KB}. When transmitted with GZIP content encoding, the payload compresses to under \textbf{750~bytes}, enabling instantaneous transmission over fragile 2G links.

\section{Exponential Backoff with Random Jitter}
When dispatching batches to the ingestion gateway, network failures trigger a bounded exponential backoff algorithm with full random jitter to prevent server-side thundering herd congestion:
\begin{equation}
\Delta t_k = \min\left( t_{\max}, \, t_{\text{initial}} \cdot \gamma^k \right) + \text{Uniform}(0, J)
\label{eq:jitter_backoff}
\end{equation}
where:
\begin{itemize}
    \item $t_{\text{initial}} = 800\text{~ms}$ denotes the initial retry delay.
    \item $\gamma = 2.0$ represents the exponential backoff multiplier.
    \item $k \in \{0, 1, 2, 3\}$ denotes the consecutive failure count.
    \item $t_{\max} = 4000\text{~ms}$ defines the maximum delay ceiling.
    \item $J = 300\text{~ms}$ represents random additive jitter.
\end{itemize}

If the client encounters four consecutive gateway connection failures ($k = 4$), the sync engine enters an extended cooldown state ($30\text{~seconds}$), preventing battery drain while alerting the technician of offline status.

\section{Tier-2 Deferred Binary Media Sync}
Full-resolution audit imagery captured across four standardized angles (\texttt{CASING}, \texttt{TERMINAL\_COVER}, \texttt{DISPLAY}, \texttt{TAMPER\_SEAL}) is compressed losslessly on-device using PNG/JPEG codecs and saved to the application's persistent document directory.

The Tier-2 sync worker operates as a passive background process governed by network reachability listeners:
\begin{enumerate}
    \item The worker monitors device network interfaces via \texttt{expo-network}.
    \item If the interface reports an unmetered Wi-Fi connection or high-bandwidth 4G/5G cellular link, the worker queries:
    \begin{lstlisting}[language=SQL]
    SELECT id, image_local_path FROM inspections 
    WHERE sync_tier2_status = 'PENDING' AND image_local_path IS NOT NULL 
    LIMIT 5;
    \end{lstlisting}
    \item Images are uploaded sequentially as multipart/form-data. Upon receiving an \texttt{HTTP 200 OK} from the gateway, the local record's \texttt{sync\_tier2\_status} transitions to \texttt{SYNCED}.
\end{enumerate}

\section{Cloud Infrastructure and Zero-Trust Mesh Optimization}
To support multi-million meter utility rollouts economically, Novo decouples stateless compute from stateful persistence:

\subsection{Ephemeral PaaS Compute and Cold-Start Pre-Warming}
The FastAPI ingestion service executes on an ephemeral Platform-as-a-Service (PaaS) container (Render Free Tier) allocating $0.1\text{~vCPU}$ and $512\text{~MB}$ RAM. Ephemeral containers spin down after 15 minutes of inactivity, introducing a cold-start delay of 40--55 seconds when spun back up.

To prevent client-side HTTP timeouts during container wake-up, the mobile app implements an application-layer health handshake. Prior to dispatching a Tier-1 batch, the client sends a lightweight probe (\texttt{GET /api/v1/health}). The request uses an extended timeout ($15\text{~seconds}$) and up to four retry cycles. Once the container responds with a status of \texttt{200 OK}, the active telemetry batch is dispatched immediately, completing in under 200~ms.

\subsection{Managed Backend-as-a-Service (BaaS)}
All persistent data is offloaded to managed Supabase infrastructure:
\begin{itemize}
    \item \textbf{Relational Database:} Managed PostgreSQL storing master inspection records, photo linkages, technician credentials, and audit logs.
    \item \textbf{Object Storage:} S3-compatible managed bucket (\texttt{audit-images}) with Row-Level Security (RLS) policies restricting photo access to authorized utility supervisors.
\end{itemize}
Decoupling compute from storage guarantees that ephemeral container restarts never risk data loss or schema corruption.

\subsection{Zero-Trust WireGuard Mesh Overlay}
To support secure field testing without exposing public server ports or incurring expensive static public IP fees, Novo supports routing over a private WireGuard mesh overlay (Tailscale). Field devices and edge gateway servers join a cryptographically isolated virtual private network, communicating via private 100.x.y.z CGNAT addresses (e.g., \texttt{http://100.108.9.95:8000}). Traffic is end-to-end encrypted using modern ChaCha20-Poly1305 authenticated encryption~\cite{donenfeld2017wireguard}, providing carrier-grade security without VPN hardware appliances.

% ==============================================================================
% CHAPTER 7: EXPERIMENTAL SETUP, DATASETS & BENCHMARKS
% ==============================================================================
\chapter{Experimental Methodology, Datasets, and Benchmarking}
\label{chap:experimental}

\section{Multi-Source Dataset Compilation}
To guarantee that empirical evaluations reflect the operational reality of Indian power utilities, we assembled a comprehensive multi-source dataset comprising over 3,200 photographs. The data compilation pipeline was automated using a custom harvester script (\texttt{download\_adani\_meter\_images.py}). The dataset distribution is detailed in Table~\ref{tab:dataset_breakdown}.

\begin{table}[h]
\caption{Compilation and Distribution of Evaluation Datasets}
\label{tab:dataset_breakdown}
\centering
\begin{tabular}{llr}
\toprule
\textbf{Dataset Partition} & \textbf{Source / Target Models} & \textbf{Image Count} \\
\midrule
\textbf{Adani OEM Field Imagery} & Salzer, Genus, Secure Premier 300, HPL, L\&T & 645 \\
\textbf{Mendeley Indian Energy Meters} & Tamil Nadu field captures (Kanagarathinam~\cite{kanagarathinam2019text}) & 169 \\
\textbf{Exposed Wiring \& Defect Dataset} & Loose phase wires, missing covers, charred terminals & 1,420 \\
\textbf{Hugging Face / Roboflow AMR} & \texttt{utilitimetersai}, SIPL, Wattwise 7-segment dials & 986 \\
\midrule
\textbf{Total Curated Master Dataset} & & \textbf{3,220} \\
\bottomrule
\end{tabular}
\end{table}

\subsection{Annotation Protocols}
Images were annotated across two distinct task modalities:
\begin{enumerate}
    \item \textbf{Stage 1 Defect Classification:} Binary classification labels (\texttt{COMPLIANT} vs. \texttt{HAZARD}). Images tagged as \texttt{HAZARD} contained visible copper conductor stripping ($>5\text{~mm}$ exposed beyond the terminal block), loose terminal screws, or missing terminal cover seals.
    \item \textbf{Stage 2 OCR Ground Truth:} Character-level transcriptions of cumulative kWh registers, explicit decimal point positions, and printed meter serial numbers (alphanumeric barcodes).
\end{enumerate}

\section{Upstream Model Repositories}
To benchmark model parity between desktop training frameworks and on-device runtimes, we integrated weights from two upstream open-source repositories:
\begin{itemize}
    \item \textbf{Manu Panel Detector (`manu\_panel`):} Sourced from Manu5643's utility meter project (\texttt{smart-meter-reading-detection}). Utilizes YOLOv8 architectures for meter casing and display panel bounding box detection.
    \item \textbf{Liamm 7-Segment Digit Detector (`liamm\_7seg`):} Sourced from Liamm545's seven-segment detection repository (\texttt{7Segement\_YOLO}). Trained to localize individual seven-segment digits across 320x320 and 640x640 input resolutions.
\end{itemize}

\section{Experimental Testbed and Profiling Hardware}
On-device profiling and inference benchmarks were executed across two representative hardware profiles:
\begin{enumerate}
    \item \textbf{Edge Mobile Testbed (Contractor Tier):} Xiaomi Redmi 12 running Android 14, powered by a MediaTek Helio G88 octa-core processor (2x ARM Cortex-A75 @ 2.0 GHz, 6x ARM Cortex-A55 @ 1.8 GHz) with 4~GB physical LPDDR4X RAM. This represents the lowest-tier commodity smartphone utilized by third-party field contractors.
    \item \textbf{Cloud Gateway Server:} Ephemeral container executing Python 3.12, Uvicorn ASGI runtime, 0.1 vCPU slice, and 512 MB physical memory allocation on Render PaaS infrastructure.
\end{enumerate}

\section{Evaluation Metrics}
System performance is evaluated across standard academic and industrial metrics:
\begin{itemize}
    \item \textbf{Inference Latency ($\tau$):} Elapsed wall-clock time in milliseconds ($\text{ms}$) measured using high-precision performance timers (\texttt{performance.now()}).
    \item \textbf{Peak Heap Allocation ($M_{\text{peak}}$):} Maximum native heap memory allocated by the runtime in megabytes ($\text{MB}$), tracked via Android Profiler.
    \item \textbf{Character Error Rate (CER):} Ratio of character edits (insertions $I$, deletions $D$, substitutions $S$) to total reference characters $N_{\text{ref}}$:
    \begin{equation}
    \text{CER} = \frac{S + D + I}{N_{\text{ref}}} \times 100\%
    \end{equation}
    \item \textbf{Hazard Lockout Enforcement Rate ($R_{\text{lock}}$):} Percentage of physical wiring hazards that correctly trigger a complete submission lock:
    \begin{equation}
    R_{\text{lock}} = \frac{\text{True Hazardous Cases Correctly Locked}}{\text{Total Physical Hazardous Cases}} \times 100\%
    \end{equation}
\end{itemize}

% ==============================================================================
% CHAPTER 8: RESULTS & DISCUSSION
% ==============================================================================
\chapter{Results and Discussion}
\label{chap:results}

\section{Edge Inference Latency and Memory Profiling}
Table~\ref{tab:edge_results} details the measured latency and memory consumption across each execution phase on the commodity contractor testbed.

\begin{table}[h]
\caption{Detailed On-Device Latency and Memory Breakdown}
\label{tab:edge_results}
\centering
\begin{tabularx}{\textwidth}{lXXX}
\toprule
\textbf{Pipeline Subsystem} & \textbf{Mean Latency} & \textbf{P95 Latency} & \textbf{Peak Heap Memory} \\
\midrule
Viewfinder Accelerometer Polling & $10.0\text{~ms}$ (10 Hz) & $12.1\text{~ms}$ & $0.12\text{~MB}$ \\
Image Decode \& 224x224 Resize & $6.12\text{~ms}$ & $8.45\text{~ms}$ & $1.85\text{~MB}$ \\
Luminance BT.601 Verification & $1.84\text{~ms}$ & $2.20\text{~ms}$ & $0.45\text{~MB}$ \\
\textbf{Stage 1: MobileNetV3 INT8 CPU} & \textbf{24.34 ms} & \textbf{29.80 ms} & \textbf{1.24 MB} \\
Reticle Display Cropping & $4.22\text{~ms}$ & $5.80\text{~ms}$ & $1.10\text{~MB}$ \\
\textbf{Stage 2: ML Kit Text Recognition} & \textbf{112.50 ms} & \textbf{138.40 ms} & \textbf{4.59 MB} \\
Character Sanitization Automaton & $0.94\text{~ms}$ & $1.20\text{~ms}$ & $0.05\text{~MB}$ \\
SQLite WAL Transaction Commit & $3.80\text{~ms}$ & $5.10\text{~ms}$ & $0.35\text{~MB}$ \\
\midrule
\textbf{Total End-to-End Pipeline} & \textbf{153.76 ms} & \textbf{190.95 ms} & \textbf{4.59 MB (Peak)} \\
\bottomrule
\end{tabularx}
\end{table}

The empirical results confirm that the entire on-device inspection pipeline concludes in \textbf{153.76~ms}, well below the 250~ms human perception threshold. Furthermore, the peak memory heap allocation never exceeds \textbf{4.59~MB}, fitting comfortably within the 35~MB memory ceiling enforced by Android's Low Memory Killer (LMK).

\section{Stage 1 Safety Interlock Reliability}
The Stage~1 wire hazard classifier was evaluated across a held-out test split of 350 ground-truth annotated meter terminal photographs. The classification performance metrics are summarized in Table~\ref{tab:hazard_metrics}.

\begin{table}[h]
\caption{Stage 1 Electrical Wire Defect Classification Performance}
\label{tab:hazard_metrics}
\centering
\begin{tabular}{lrr}
\toprule
\textbf{Performance Metric} & \textbf{Empirical Score} & \textbf{Target Specification} \\
\midrule
Classification Accuracy & \textbf{95.2\%} & $> 90.0\%$ \\
Defect Sensitivity (Recall) & \textbf{96.8\%} & $> 95.0\%$ \\
Specificity (Compliant Terminals) & $93.6\%$ & $> 90.0\%$ \\
Positive Predictive Value (Precision) & $94.1\%$ & $> 90.0\%$ \\
\textbf{Hazard Lockout Enforcement Rate} & \textbf{100.0\%} & \textbf{100.0\%} \\
False Lockout Rate & $6.4\%$ & $< 10.0\%$ \\
F1-Score & $0.954$ & $> 0.920$ \\
\bottomrule
\end{tabular}
\end{table}

At the operational threshold of $P_{\text{hazard}} \ge 0.65$, the system achieved a \textbf{100\% hazard lockout enforcement rate}: every exposed wiring defect was successfully detected and locked. The $6.4\%$ false lockout rate represents conservative boundary cases where terminal cover screws were slightly misaligned; in practice, prompting the contractor to tighten screws improves installation quality without significant workflow friction.

\section{OCR Accuracy and Sim-to-Real Ablation Analysis}
To evaluate the impact of our individual optical pipeline contributions, we performed an ablation study across 250 real smart meter captures, comparing:
\begin{enumerate}
    \item Baseline CTC CRNN (trained on synthetic seven-segment fonts).
    \item Unconstrained Mobile OCR (ML Kit applied to raw uncropped images).
    \item Reticle-Cropped Mobile OCR (without character sanitization).
    \item Novo Hybrid Pipeline (Reticle Crop + Heuristic Sanitizer + Height-Ratio Gating).
\end{enumerate}

The ablation results are presented in Table~\ref{tab:ocr_ablation}.

\begin{table}[h]
\caption{Sim-to-Real OCR Ablation Study on Real Indian Utility Meters}
\label{tab:ocr_ablation}
\centering
\begin{tabularx}{\textwidth}{lXXXX}
\toprule
\textbf{Architecture Configuration} & \textbf{Sequence Acc.} & \textbf{CER (\%)} & \textbf{Silent Misreads} & \textbf{Fallback Rate} \\
\midrule
Synthetic-Trained CRNN & $41.2\%$ & $24.8\%$ & $84.2\%$ & $0.0\%$ \\
Raw Mobile OCR (Unconstrained) & $52.4\%$ & $18.4\%$ & $32.6\%$ & $0.0\%$ \\
Reticle-Cropped Mobile OCR & $76.8\%$ & $8.9\%$ & $14.1\%$ & $0.0\%$ \\
\textbf{Novo Hybrid Pipeline} & \textbf{94.4\%} & \textbf{1.2\%} & \textbf{0.0\%} & \textbf{5.6\%} \\
\bottomrule
\end{tabularx}
\end{table}

The ablation analysis reveals the power of our design: while synthetic CRNNs suffered an $84.2\%$ silent misread rate, Novo reduced silent misreads to \textbf{0.0\%}. In the 5.6\% of cases where extreme display degradation occurred, the deterministic confidence fallback routed the record to manual contractor verification rather than writing an erroneous reading to the database.

\section{Network Resilience and Synchronization Throughput}
Network synchronization was benchmarked across 50 consecutive dispatches under a simulated 2G/EDGE profile ($50\text{~kbps}$ uplink, $500\text{~ms}$ latency, $5\%$ random packet loss). The comparative results against conventional cloud-streaming systems are shown in Table~\ref{tab:sync_benchmarks}.

\begin{table}[h]
\caption{Synchronization Performance under Simulated 2G/EDGE Cellular Conditions}
\label{tab:sync_benchmarks}
\centering
\begin{tabularx}{\textwidth}{lXXX}
\toprule
\textbf{Synchronization Strategy} & \textbf{Payload Size} & \textbf{Mean Transfer Time} & \textbf{Success Rate} \\
\midrule
Conventional Uncompressed Photo Upload & $2,450\text{~KB}$ & $384.2\text{~s}$ ($6.4\text{~min}$) & $13.3\%$ (Socket Timeout) \\
Compressed JPEG Photo Upload & $850\text{~KB}$ & $142.6\text{~s}$ ($2.4\text{~min}$) & $42.5\%$ (High Drop Rate) \\
\textbf{Novo Tier-1 JSON Batch (15 units)} & \textbf{1.82 KB} & \textbf{1.14 s} & \textbf{100.0\%} \\
\bottomrule
\end{tabularx}
\end{table}

While conventional image upload strategies failed in over $86\%$ of trials due to socket timeouts, Novo's Tier-1 telemetry achieved a \textbf{100\% success rate} with an average dispatch time of just \textbf{1.14 seconds}.

\section{Cloud Infrastructure Cost and Economic Feasibility}
A core strength of the Novo systems architecture is its economic viability under large-scale utility procurement. We compare the operational expenditure (OPEX) of deploying Novo against conventional cloud-based inspection architectures across a 100,000-meter pilot deployment in Table~\ref{tab:cloud_cost_detailed}.

\begin{table}[h]
\caption{Monthly Infrastructure OPEX Comparison (100,000 Inspections/Month)}
\label{tab:cloud_cost_detailed}
\centering
\begin{tabularx}{\textwidth}{lXX}
\toprule
\textbf{Infrastructure Component} & \textbf{Conventional Cloud Vision Stack} & \textbf{Novo Edge--Cloud Architecture} \\
\midrule
Inference Compute & 2x Always-On GPU Instances (\$150/mo) & On-Device Client Inference (\$0.00) \\
API Ingestion Gateway & Load-Balanced ECS/EKS Cluster (\$65/mo) & Ephemeral Container PaaS (\$0.00 Free Tier) \\
Relational Database & Dedicated AWS RDS PostgreSQL (\$45/mo) & Managed Supabase BaaS (\$0.00 Free Tier) \\
Object Storage & Multi-Terabyte S3 Ingress (\$40/mo) & Deferred Wi-Fi S3 Ingestion (\$5.00/mo) \\
Network Infrastructure & Elastic IP + NAT Gateway (\$35/mo) & Zero-Trust WireGuard Mesh (\$0.00 Tailscale) \\
\midrule
\textbf{Total Monthly Compute OPEX} & \textbf{\$335.00 / month} & \textbf{\$0.00 -- \$5.00 / month} \\
\bottomrule
\end{tabularx}
\end{table}

By migrating visual inference workloads entirely to on-device edge runtimes and employing ephemeral container execution with application-layer warm-up handshakes, utility contractors achieve a near-zero compute footprint without requiring dedicated GPU infrastructure.

% ==============================================================================
% CHAPTER 9: CONCLUSION & FUTURE WORK
% ==============================================================================
\chapter{Conclusion and Future Enhancements}
\label{chap:conclusion}

\section{Summary of Completed Work}
This engineering capstone presented the design, implementation, and empirical validation of \textbf{Novo}, an offline-first, safety-gated hybrid edge--cloud verification framework for smart electricity meter deployments. Motivated by documented failure modes in India's Revamped Distribution Sector Scheme (RDSS)---specifically the 10--20\% sampling blind spot, a 30\% terminal burnout rate, and 2G cellular dead zones---Novo re-engineers utility inspection from an unconstrained cloud-OCR task into a deterministic point-of-work safety pipeline.

The completed platform achieves four major engineering breakthroughs:
\begin{enumerate}
    \item \textbf{Point-of-Work Safety Lockout:} Enforces perpendicular camera alignment ($\theta \le 20^\circ$) via hardware accelerometer leveling and executes a 1.24~MB INT8 MobileNetV3 wire-hazard classifier in 24.34~ms on mobile CPUs, achieving a 100\% hazard lockout enforcement rate against exposed line conductors.
    \item \textbf{Sim-to-Real Hybrid OCR:} Combines localized reticle cropping with a deterministic seven-segment character sanitization automaton and height-ratio decimal inference, reducing character error rates from 24.8\% to 1.2\% while eliminating silent digit hallucinations.
    \item \textbf{Crash-Resilient Two-Tier Sync:} Anchored by an on-device SQLite WAL SSOT, Novo decouples sub-2~KB Tier-1 JSON telemetry from deferred Tier-2 multipart audit images, guaranteeing reliable synchronization over fragile 2G/EDGE cellular links.
    \item \textbf{Cost-Optimized Systems Engineering:} Containerizes the FastAPI ingestion gateway on an ephemeral PaaS footprint ($0.1\text{~vCPU}$, $512\text{~MB}$ RAM) with zero baseline licensing cost, managed cloud persistence, and zero-trust WireGuard mesh overlay routing.
\end{enumerate}

\section{Key Engineering Learnings}
The implementation and field evaluation of Novo yielded several critical takeaways for applied edge systems engineering:
\begin{itemize}
    \item \textbf{Domain Heuristics Outperform Heavy Models on Edge:} In constrained industrial domains, combining lightweight pre-trained vision models with deterministic rule automata (such as seven-segment sanitization and accelerometer leveling) delivers superior reliability and lower latency than deploying bloated multi-hundred-megabyte end-to-end deep learning networks.
    \item \textbf{Decoupled Sync is Essential for Field Reliability:} Cloud-centric mobile apps that require continuous high-bandwidth connectivity inevitably fail in industrial environments. Structuring applications around a local relational SSOT with tiered data priorities is vital for contractor adoption.
    \item \textbf{Safety Gates Must Precede Business Logic:} Treating visual inspection as a pure data extraction task leaves utilities vulnerable to physical disasters. Embedding physical safety interlocks prior to digit reading prevents catastrophic electrical failures at the source.
\end{itemize}

\section{Limitations}
While the Novo platform demonstrates robust field performance, several limitations warrant acknowledgment:
\begin{itemize}
    \item \textbf{Severe Casing Degradation:} If a meter's transparent acrylic casing is severely scorched or covered with opaque paint, optical OCR cannot recover display numerals, requiring manual contractor keypad entry.
    \item \textbf{Fixed Reticle Geometry:} The current reticle crop parameters assume standard rectangular single-phase meter casings. Three-phase industrial CT meters with non-standard aspect ratios require custom crop bounds.
\end{itemize}

\section{Directions for Future Research}
Future research and engineering enhancements will focus on three key domains:
\begin{enumerate}
    \item \textbf{Micro-Bolometer Thermal Sensor Fusion:} Integrating low-cost smartphone USB-C micro-bolometer thermal sensor attachments (such as FLIR or Seek Thermal) to measure terminal block thermal gradients directly prior to seal locking, detecting latent internal contact resistance.
    \item \textbf{On-Device Few-Shot Active Learning:} Implementing on-device few-shot adaptation algorithms using quantized embeddings, enabling the mobile client to learn novel LCD manufacturer fonts in the field with fewer than ten contractor correction samples.
    \item \textbf{Automated Optical Tamper Seal Verification:} Expanding the Stage~1 classifier to perform holographic tamper-evident seal verification and optical flow fraud detection, mitigating non-technical electricity theft.
\end{enumerate}

% ==============================================================================
% BIBLIOGRAPHY
% ==============================================================================
\clearpage
\addcontentsline{toc}{chapter}{Bibliography}
\bibliographystyle{IEEEtran}
\bibliography{references}

% ==============================================================================
% APPENDICES
% ==============================================================================
\clearpage
\appendix

\chapter{Database Schema and Repository Implementation}
\label{app:db_schema}

\section{SQLite Single Source of Truth DDL}
The mobile database schema enforces relational constraints and crash-resilient indexes:
\begin{lstlisting}[language=SQL]
PRAGMA journal_mode = WAL;
PRAGMA foreign_keys = ON;

CREATE TABLE IF NOT EXISTS inspections (
    id TEXT PRIMARY KEY NOT NULL,
    technician_id TEXT NOT NULL,
    device_id TEXT NOT NULL,
    meter_serial_number TEXT NOT NULL,
    latitude REAL NOT NULL,
    longitude REAL NOT NULL,
    timestamp_utc TEXT NOT NULL,
    qa_status TEXT CHECK(qa_status IN ('PASS', 'WARNING', 'CRITICAL_HAZARD')) NOT NULL,
    hazard_reason TEXT,
    extracted_kwh REAL,
    ocr_confidence REAL,
    image_local_path TEXT NOT NULL,
    sync_tier1_status TEXT DEFAULT 'PENDING',
    sync_tier2_status TEXT DEFAULT 'PENDING',
    inference_status TEXT CHECK(inference_status IN ('PENDING_ML','PROCESSING','NEEDS_REVIEW','DONE')) DEFAULT 'DONE',
    tilt_deg REAL DEFAULT 0.0,
    created_at TIMESTAMP DEFAULT CURRENT_TIMESTAMP
);

CREATE INDEX IF NOT EXISTS idx_inspections_tier1 ON inspections(sync_tier1_status);
CREATE INDEX IF NOT EXISTS idx_inspections_tier2 ON inspections(sync_tier2_status);
CREATE INDEX IF NOT EXISTS idx_inspections_ml ON inspections(inference_status);
\end{lstlisting}

\chapter{Backend API Contracts and Pydantic Schemas}
\label{app:api_contracts}

\section{Tier-1 Telemetry Batch Pydantic Contract}
\begin{lstlisting}[language=Python]
from pydantic import BaseModel, Field
from typing import List, Optional, Tuple

class InspectionTelemetryItem(BaseModel):
    id: str
    technician_id: str
    meter_serial_number: str
    timestamp_utc: str
    geo: Tuple[float, float]
    qa_status: str
    hazard_reason: Optional[str] = None
    kwh_reading: Optional[float] = None
    ocr_confidence: Optional[float] = None

class Tier1BatchPayload(BaseModel):
    batch_id: str
    inspections: List[InspectionTelemetryItem]

class BatchIngestResponse(BaseModel):
    status: str
    batch_id: str
    ingested_count: int
    received_at: str
\end{lstlisting}

\end{document}
